\chapter{Ausgewählte Konfigurations- und Code-Auszüge}
\label{app:listings}

Dieser Anhang sammelt drei Auszüge aus den in Kapitel~\ref{cha:realisierung} beschriebenen Skripten und Konfigurationsdateien.
Sie dienen als Beleg für die im Fliesstext zusammengefasste Funktionsweise und erlauben es, die Pipeline anhand der konkreten Dateiformate nachzuvollziehen.

\section{Schedule-Datei}
\label{app:listing_schedule}

Die von \texttt{generate-wf-config.py} erzeugte Datei \texttt{schedule.json} dient als Ground Truth für die spätere Zuordnung von Pcap-Abschnitten zu Seitenaufrufen (vgl.\ Abschnitt~\ref{sec:realisierung_generate_wf_config}).
Pro Monitor enthält sie eine Liste der geplanten Besuche mit Startzeit, Zielport und URL.

\begin{lstlisting}[
    basicstyle=\ttfamily\footnotesize,
    breaklines=true,
    frame=single,
    captionpos=b,
    label={lst:schedulejson},
    caption={Auszug aus \texttt{schedule.json}: für jeden Monitor werden Startzeit, Zielport und URL der geplanten Besuche festgehalten.}
]
"monitor0": [
    {
        "start_time": 1200,
        "port": 8000,
        "url": "http://129.114.108.192:8000/Eurovision_Song_Contest_2018"
    },
    {
        "start_time": 1230,
        "port": 8010,
        "url": "http://129.114.108.192:8010/Ionizing_radiation"
    }
]
\end{lstlisting}

\section{URL-Liste}
\label{app:listing_urls}

Die von \texttt{generate-urls.py} erzeugte Datei \texttt{urls.txt} legt fest, welche Wikipedia-Seiten den Monitor-Hosts zugewiesen werden (vgl.\ Abschnitt~\ref{sec:realisierung_webserver}).
Pro Zeile sind die IP des Webservers, der für die Seite verwendete Port und die vollständige URL angegeben.

\begin{lstlisting}[
    basicstyle=\ttfamily\footnotesize,
    breaklines=true,
    frame=single,
    captionpos=b,
    label={lst:urltxt},
    caption={Beispielhafter Auszug aus \texttt{urls.txt}. Jede Zeile enthält IP, Port und URL einer ausgewählten Wikipedia-Seite.}
]
129.114.108.192 8000 http://129.114.108.192:8000/Eurovision_Song_Contest_2018
129.114.108.192 8001 http://129.114.108.192:8001/Andhra_Pradesh
129.114.108.192 8002 http://129.114.108.192:8002/Rajon_Rondo
129.114.108.192 8003 http://129.114.108.192:8003/Hagerman,_Idaho
129.114.108.192 8004 http://129.114.108.192:8004/Echzell
\end{lstlisting}

\section{NEWNYM-Skript}
\label{app:listing_newnym}

Das Skript \texttt{newnym.py} erzwingt vor jedem Seitenbesuch einen frischen Tor-Circuit (vgl.\ Abschnitt~\ref{sec:realisierung_newnym}).
Es öffnet einen Socket zum lokalen Tor-Control-Port und sendet das Signal \texttt{NEWNYM}.

\begin{lstlisting}[
    basicstyle=\ttfamily\footnotesize,
    frame=single,
    captionpos=b,
    label={lst:newnym},
    caption={\texttt{newnym.py}: minimaler Tor-Control-Port-Client, der vor jedem Besuch einen frischen Circuit erzwingt.}
]
import socket

with socket.socket(socket.AF_INET, socket.SOCK_STREAM) as s:
    s.connect(("127.0.0.1", 9051))
    s.sendall(b"AUTHENTICATE\r\n"); s.recv(1024)
    s.sendall(b"SIGNAL NEWNYM\r\n"); s.recv(1024)
\end{lstlisting}
